\documentclass[10pt]{article}
\usepackage[margin=20mm]{geometry}
\usepackage{amsmath,amssymb,fontspec,unicode-math,xurl,hyperref}
\IfFontExistsTF{DejaVu Serif}{\setmainfont{DejaVu Serif}}{\setmainfont{Latin Modern Roman}}
\setmathfont{latinmodern-math.otf}
\setlength{\parindent}{0pt}\setlength{\parskip}{6pt}\sloppy
\begin{document}
Abstract: The hierarchical matching approach has a plausible repair path preserving its headline asymptotic tradeoff, but the published augmentation construction, its local complexity, and several update instructions require correction. The released implementation supplies useful counterevidence to a claim of a broken matching engine: it uses directed residual transitions and records actual point distances. Its deletion timing, however, is stale. The augmentation specification and the update invariant are Major defects in the published proof of the main dynamic theorem. The crossing calculation and released timing instrumentation remain bounded Minor findings. A repair may preserve the headline theorem, but it is not supplied by the published argument. The lower bound must state the representation and recourse model it needs. All 25 physical PDF pages and all five original figures were inspected, together with the complete small source release. Only reviewer-written finite checks were executed. Scientific adjudication remains pending.\par
Author: AI-assisted technical review\par
Date:\par
Title: Technical Verification of Fully Dynamic Euclidean Bi-Chromatic Matching in Sublinear Update Time\par
\# Identity, sources, and coverage\par
The reviewed paper is Gramoz Goranci, Peter Kiss, Neel Patel, Martin P. Seybold, Eva Szilagyi, and David Zheng, Fully Dynamic Euclidean Bi-Chromatic Matching in Sublinear Update Time, ICML 2025, I25-43715 (0346) [Retained latex: \textbackslash{}cite\{paper\}]. The frozen PDF SHA-256 is\par
.\par
The complete 25-page paper, including appendices A–F, was read. Locations below are physical PDF pages. Figures 1–5 and the relevant proof and pseudocode pages were actually viewed. All 14 acquired packet files passed their frozen manifest hashes.\par
The author repository [Retained latex: \textbackslash{}cite\{code\}] was read at commit , dated 21 November 2025: all 747 lines of , the three input-generator scripts, README, build script and CMake file. This later public snapshot is not assumed to be the version that produced the paper’s experiments. No author program, build, benchmark, or data generator was run.\par
\# The algorithmic argument and checks that survive\par
Each cell matches a balanced subset of the points forwarded by its children and forwards a monochromatic excess set. At leaves the matching is exact. Internal cells aggregate excess points to at most \(p^2\) child centers, solve transportation there, and expand the assignment to actual points. The random-shift analysis charges cross-cell edges of an optimal matching at each level; \(\log_p U=O(1/\epsilon)\) for polynomial spread and the chosen \(p\).\par
The static theorem’s high-probability wording is not, by itself, an error. Appendix C, page 14, explicitly repeats independent shifts \(O(\log n)\) times and selects a minimum-cost solution. The exponent slack in \(p=n^{\epsilon/10}\) can absorb this amplification for fixed \(\epsilon\). The dynamic theorem states an expectation. A guarantee against an adaptive adversary that observes the shift or output is not established by the fixed-instance crossing calculation; this scope should be explicit.\par
Claims E.1–E.2 contain the standard optimal-matching extension idea. For an optimal old matching, a shortest alternating augmenting path produces an optimal extension: the other components of the symmetric difference are nonnegative-cost alternating cycles. Deleting matched pairs leaves an optimal restriction on the remaining endpoints, so the proposed removal-and-reinsertion reduction is reasonable when its distinct-endpoint cases are handled. Independent enumeration of 100 instances with zero to four old pairs confirmed the cost identity for the corrected directed residual network, to numerical error below \(10^{-9}\). These small tests support that repair; they do not prove all dynamic invariants or runtime bounds.\par
\# Confirmed corrections\par
Problem 1 (The written augmentation graph and its complexity are incorrect).  Appendix E.1, page 21 (the lemma heading prints A.4, while the surrounding text and page 22 refer to E.4), does not specify a valid residual network as written, and its \(O(p^3)\) bound in Lemma E.4 does not follow from the graph it constructs.\par
First, the new vertices \(x',y'\) are connected only to old opposite-color centers. Their direct edge is missing. With an empty old matching, two newly inserted opposite-color points must be matched to each other, but the specified graph has no \(x'\)–\(y'\) path. This case is allowed by the lemma and arises naturally when a cell first obtains both colors.\par
Second, a single undirected signed edge per pair of centers does not enforce alternating paths. A transportation residual network needs forward red-to-blue arcs of positive distance and reverse blue-to-red arcs, of negative distance, wherever positive flow can be removed. The forward capacity is not lost merely because the same center pair already carries some flow. The text instead assigns either a positive or negative weight according to whether the pair is used. Interpreting an undirected negative edge as two directed arcs also creates a negative two-cycle. The endpoint paragraph additionally copies \(x'\) where \(y'\) is intended.\par
Third, with \(k=p^2\) centers and half of each color, the constructed complete bipartite graph has \[|V|=\Theta(p^2),\qquad |E|=\Theta(p^4).\] Explicit construction alone can take \(\Theta(p^4)\), and the quoted Bellman–Ford bound is \(O(p^6)\), not \(O(p^3)\). For example, \(p=16\) gives 256 centers and 16,384 old bipartite edges, already exceeding \(p^3=4096\).\par
These defects are Major because the proof of Theorem 1.1 on physical page 22 explicitly invokes this augmentation routine for its update time and correctness. The empty-state example invalidates the specified subroutine on an admitted state, and the stated graph counts invalidate its cited runtime derivation. The intended directed residual algorithm is a credible repair, and the headline exponent might survive: \(p^6=n^{0.6\epsilon}\) for the stated \(p\), while even an \(O(p^2)\)-node rebuild has exponent below the reserved \(p^{10}\) scale. That slack does not validate the present lemma. A replacement network, proof of its augmentation property, and complete update/rebuild accounting are required. Neither the existence of an algorithm with the headline tradeoff nor the inspected implementation is refuted by this finding.\par
The inspected C++ is counterevidence to a broader accusation. Lines 313–350 use directed positive forward transitions and negative backward transitions along matched edges. An opposite-color exposed point is a candidate destination, including the direct new-pair case. Lines 357–408 recover the alternating path. Thus the literal missing-edge/direction defect in the paper is not established in this implementation.\par
Problem 2 (The printed update pseudocode violates its excess invariant).  Algorithms 1, 3 and 5 on pages 18–19 and 23 require coordinated corrections, rather than literal implementation.\par
For Algorithm 5, take a non-marginal leaf containing one blue point and no red points. Insert one red point. Line 41 recomputes the imbalance as zero; the \(\eta\geq0\) branch then adds the red point to the old blue excess set without matching them. A balanced leaf is left with a bichromatic excess set and no match. This contradicts the invariant used in Appendix D.2. If the intended branch uses the old imbalance, the other branch still uses \(E_v\leftarrow E_v\cup\{b\}\) at line 47, where removal is needed after matching \(b\).\par
Algorithm 1 decreases the imbalance after inserting red into a blue-excess cell, although \(|A|-|B|\) increases by one. Algorithm 3 recursively calls Insert-Update during deletion. Its leaf case can add the deleted point itself to the excess set. These are explicit specification failures. The repair is to state pre/post-update imbalance consistently, remove points when they are consumed, forward only remaining endpoints, and recurse with the correct operation. The release’s , lines 511–552, follows endpoint changes and calls removal functions instead; no identical implementation failure is inferred. The finding is Major for the published algorithm: the correctness proof of Theorem 1.1 on page 22 assumes precisely the post-update monochromatic-excess property violated by this trace. A local example can be embedded in a globally balanced pair update by placing the other new point in a separate cell. The invariant prose and later code are useful evidence for an intended repair, but they do not prove that the printed procedures preserve the claimed state. Corrected procedures and a case-complete invariant proof are required; this does not establish failure of the later implementation.\par
Problem 3 (The crossing estimate has the wrong scale factor).  Claim C.6, page 16, divides expected crossings by cell width; the subsequent cost proof needs multiplication.\par
For a fixed horizontal edge of length \(L\leq\lambda\) under a uniform grid shift of width \(\lambda\), the crossing probability is exactly \(L/\lambda\). Thus its contribution satisfies \[\lambda\,\mathbb E[\operatorname{Cross}]=L,
 \qquad \mathbb E[\operatorname{Cross}]/\lambda=L/\lambda^2.\] The correct upper bound is \(\lambda_\ell\mathbb E[\operatorname{Cross}(\mu^*,\ell)]=O(c(\mu^*))\), using an upper crossing-probability bound. The printed proof instead writes an \(\Omega(L/\lambda)\) probability and says it contributes \(L\). Long edges also require truncating a probability bound at one. Equation (5), page 17, is consistent with the corrected identity.\par
The bad-point counting in Claim C.3 also needs care: if half of \(B_\ell\) bad points are incident to crossing edges, an edge may have two such endpoints, giving \(B_\ell\leq4\operatorname{Cross}\) from that argument. Each swap eliminates at least two bad points and costs at most \(4\lambda_\ell\), which recovers the stated constant 8. The displayed \(B_\ell\leq\operatorname{Cross}/2\) on page 16 has its relationship reversed. These are Minor local proof errors: the exact uniform-shift upper bound and endpoint accounting above establish the inequality actually used in equation (5), under the same fixed-instance assumptions, with the same expected \(O(1/\epsilon)\) approximation bound. The classification rests on this verified surviving argument, not on the ease of editing the formulas.\par
Problem 4 (Deletion rows in the released logger reuse an old runtime).  The public snapshot does not measure deletion update time in its reported column.\par
Insertion assigns inside lines 561–565. Deletion, lines 570–576, never assigns it. The per-operation reset at lines 64–67 resets only changed-edge and iteration counters, and each operation then prints the statistics at line 723. Symbolically, an insertion assigning time 123 followed by any deletion produces another row with time 123, regardless of the deletion duration. If no insertion has set it, the initial value is \(-1\). Also, insertion’s timer starts after point storage and tree traversal in , which is called at line 557 before the timer at line 561, so it measures the recursive matching modification, not the complete exposed insertion operation.\par
This is a confirmed Minor measurement defect in the release, relevant to reconstructing the fully dynamic benchmark of Figure 5. It does not establish that the historical figure used this logger or has a particular error size. Full-operation timing should bracket both insertion and deletion consistently, and the historical experimental version and raw logs should be supplied. The exact-baseline branch at lines 711–713 is only a TODO in this snapshot; that is a reproduction limit, not evidence that the paper’s separately reported exact baseline was never computed.\par
\# Scope questions and empirical limits\par
The lower bound in Theorem F.1, pages 24–25, converts \(o(n)\) update time into \(o(n)\) changed edges on some update. That conversion needs an explicit-maintenance or charged-recourse model. On its alternating-line construction, the exact matching can instead be represented by the left endpoint, right endpoint, and the rule “pair consecutive positions from the left.” Each insertion changes two stored endpoints; the exact cost is \(n+i\), and enumerating the full matching later costs \(O(n+i)\). The independent check confirms this rule for the displayed sequence. This does not disprove a general two-dimensional lower bound. It shows why the given proof does not cover unrestricted implicit representations or cost-only queries. Under explicit edge maintenance, the recourse argument is persuasive; the theorem and introductory claim should state that restriction. This is recorded as a model-scope question rather than an unconditional refutation.\par
Figures 2–5 were checked against their captions and prose. Figure 3 includes exact comparisons only through roughly 2000 samples; Figure 4 extends approximate curves to one million but has no exact comparator. Consequently the less-than-two empirical approximation claim is directly supported only where an exact comparator is available, not certified by visual convergence at a million points. Figure 5 depicts distribution drift and update times, but no repeated-run uncertainty is supplied. The source packet lacks the historical taxi preprocessing, raw benchmark outputs, complete exact-baseline implementation and plotting pipeline. These numerical results were not independently reproduced.\par
The program’s point manager adds and subtracts actual point distances when matching and unmatching (138–156). Thus the basic algorithm’s center-cost estimate should not be confused with the implementation’s reported total matching cost. The basic prose explicitly describes an estimate on page 6; equality between center cost and expanded matching cost is not assumed in this review.\par
\# Disposition and evidence\par
Recommendation: Major revision, driven by the invalid augmentation specification and post-update invariant on which the main dynamic theorem relies. The crossing calculation, timing instrumentation, and theorem-scope qualifications also require the bounded corrections stated above. No quota was applied. No conclusion about historical experiment contamination, universal correctness of the implementation, or an adaptive-adversary guarantee is made.\par
The companion evidence includes physical-page text, source hashes, original-page render hashes, reviewer-written and , claim-level findings and coverage, isolated compilation logs, and actual final-PDF inspection records. Task , token ; selected model , effort . The predecessor is preserved. Focused independent AI verification corrected severity using the actual theorem dependencies and checked the counterevidence; it does not certify a replacement algorithm or historical experiments. Scientific adjudication is pending and human reviews are empty.\par
2 G. Goranci et al., Fully Dynamic Euclidean Bi-Chromatic Matching in Sublinear Update Time, ICML 2025. https://proceedings.mlr.press/v267/goranci25a.html. Author repository, pinned source snapshot. https://github.com/Zhengdw/dyn-euc-match/tree/6fa0ce9e2eaac5187a513f165c2959cc0b6a288a.\par
\end{document}
